\documentclass[12pt,a4paper]{article}

\usepackage[T1]{fontenc}
\usepackage[utf8]{inputenc}
\usepackage[romanian]{babel}
\usepackage{amsmath,amssymb}
\usepackage{booktabs}
\usepackage{geometry}
\usepackage{hyperref}
\geometry{margin=2.5cm}

\title{Parametrizarea temporală a traiectoriei și extinderea regulatorului PID}
\author{Raport intermediar pentru proiectul de diplomă}
\date{7 august 2026}

\begin{document}
\maketitle

\section{Scopul etapei}

Etapa descrisă în acest raport realizează trecerea de la urmărirea pur
geometrică a unei căi la urmărirea unei traiectorii parametrizate în timp.
Această modificare este necesară pentru o comparație riguroasă între PID, LQR
și MPC. Cele trei regulatoare vor primi aceeași referință, aceeași stare de
eroare, aceleași viteze de anticipare și aceleași limite ale comenzilor.

În formularea anterioară, punctul de referință era ales în funcție de proiecția
robotului pe cale. Dacă robotul era încetinit de o perturbație, referința
încetinea implicit împreună cu el. Eroarea longitudinală nu putea măsura în mod
real întârzierea față de evoluția dorită. În formularea actuală, un robot
virtual se deplasează pe cale conform unui ceas independent de robotul simulat.
Astfel, perturbațiile produc atât eroare transversală și unghiulară, cât și
eroare longitudinală măsurabilă.

\section{Referința geometrică și parametrizarea temporală}

Fie calea plană descrisă prin coordonata curbilinie $s$:
\begin{equation}
  \mathbf{q}_p(s) =
  \begin{bmatrix}x_p(s) & y_p(s) & \theta_p(s)\end{bmatrix}^{T}.
\end{equation}
Punctele citite din fișierul CSV sunt conectate prin segmente, iar lungimea
cumulată definește progresul $s\in[0,L]$. Din geometrie se determină tangenta
$\theta_p(s)$ și curbura $\kappa(s)$.

Parametrizarea temporală introduce funcția $s_r(t)$. Traiectoria robotului
virtual devine
\begin{equation}
  \mathbf{q}_r(t)=\mathbf{q}_p(s_r(t)), \qquad
  v_r(t)=\dot{s}_r(t), \qquad
  \omega_r(t)=\kappa(s_r(t))v_r(t).
\end{equation}
Referința este deci compusă din poziție, orientare, viteză liniară și viteză
unghiulară. Ea depinde numai de timpul experimentului, nu de poziția robotului
controlat.

\subsection{Construirea profilului de viteză}

Calea este reeșantionată spațial la un pas maxim de $0{,}01$ m. Pentru fiecare
nod se pornește de la viteza nominală, redusă în funcție de curbură. Se aplică
și condiția
\begin{equation}
  |\omega_r|=|\kappa v_r|\leq \omega_{r,\max}.
\end{equation}
Robotul virtual pornește și se oprește cu viteză nulă. O parcurgere înainte
limitează accelerația:
\begin{equation}
  v_i \leq \sqrt{v_{i-1}^{2}+2a_{\max}\Delta s},
\end{equation}
iar o parcurgere înapoi impune posibilitatea frânării până la capăt:
\begin{equation}
  v_{i-1} \leq \sqrt{v_i^{2}+2d_{\max}\Delta s}.
\end{equation}
Durata fiecărui interval este obținută presupunând accelerație constantă:
\begin{equation}
  \Delta t_i=\frac{2\Delta s}{v_i+v_{i+1}}.
\end{equation}
Între noduri se interpolează poziția curbilinie prin ecuațiile mișcării cu
accelerație constantă. Configurația curentă folosește viteza nominală de
$0{,}4$ m/s, accelerația de referință de $0{,}5$ m/s$^2$ și o frânare
conservatoare de $0{,}1$ m/s$^2$.

\section{Starea de eroare comună}

Poziția robotului este $\mathbf{q}=[x\;y\;\theta]^T$. Diferența dintre poziția
referinței temporale și poziția estimată este rotită în reperul mobil al
robotului:
\begin{equation}
\begin{bmatrix}e_x\\e_y\end{bmatrix}=
\begin{bmatrix}
\cos\theta & \sin\theta\\
-\sin\theta & \cos\theta
\end{bmatrix}
\begin{bmatrix}x_r-x\\y_r-y\end{bmatrix},
\qquad
e_{\theta}=\mathrm{wrap}(\theta_r-\theta).
\end{equation}
Aici $e_x$ este eroarea longitudinală, $e_y$ eroarea laterală, iar
$e_{\theta}$ eroarea de orientare. Aceasta va fi starea comună
$\mathbf{e}=[e_x\;e_y\;e_{\theta}]^T$ utilizată de PID, LQR și MPC.

Pentru modelul cinematic diferențial
\begin{equation}
  \dot{x}=v\cos\theta,\quad
  \dot{y}=v\sin\theta,\quad
  \dot{\theta}=\omega,
\end{equation}
dinamica exactă a erorilor în reperul mobil este
\begin{align}
  \dot e_x &= \omega e_y-v+v_r\cos e_{\theta},\\
  \dot e_y &= -\omega e_x+v_r\sin e_{\theta},\\
  \dot e_{\theta} &= \omega_r-\omega.
\end{align}
Această formă arată explicit că ambele intrări, $v$ și $\omega$, influențează
urmărirea traiectoriei.

În jurul traiectoriei nominale, pentru erori mici și comenzile
$v=v_r+\delta v$, $\omega=\omega_r+\delta\omega$, se obține modelul liniar
variabil în timp
\begin{equation}
  \dot{\mathbf e}=A(t)\mathbf e+B\delta\mathbf u,
\end{equation}
cu
\begin{equation}
A(t)=
\begin{bmatrix}
0 & \omega_r(t) & 0\\
-\omega_r(t) & 0 & v_r(t)\\
0 & 0 & 0
\end{bmatrix},\qquad
B=
\begin{bmatrix}
-1 & 0\\
0 & 0\\
0 & -1
\end{bmatrix},\qquad
\delta\mathbf u=
\begin{bmatrix}\delta v\\\delta\omega\end{bmatrix}.
\end{equation}
Acesta este punctul de plecare direct pentru LQR și MPC; nu este necesară
identificarea unui model cinematic deja cunoscut.

\section{Regulatorul PID extins la două comenzi}

Regulatorul PID actual furnizează aceeași pereche de corecții care va fi
furnizată ulterior de LQR și MPC:
\begin{equation}
  \delta\mathbf u_{PID}=
  \begin{bmatrix}\delta v_{PID} & \delta\omega_{PID}\end{bmatrix}^{T}.
\end{equation}

Eroarea longitudinală este tratată de un PID separat:
\begin{equation}
  \delta v_{PID}=K_{p,x}e_x+K_{i,x}\int e_xdt+K_{d,x}\dot e_x.
\end{equation}
Pentru mișcarea laterală se păstrează structura cascadată. Bucla externă
transformă $e_y$ într-o corecție limitată de orientare:
\begin{equation}
  \Delta\theta_c=\mathrm{sat}\left(PID_y(e_y),
  -\Delta\theta_{\max},\Delta\theta_{\max}\right),
\end{equation}
iar bucla internă reglează orientarea corectată:
\begin{equation}
  e_{\theta,c}=\mathrm{wrap}(\theta_r+\Delta\theta_c-\theta),
  \qquad \delta\omega_{PID}=PID_{\theta}(e_{\theta,c}).
\end{equation}
Prin această structură, o rotație bruscă a robotului este observată imediat de
bucla internă; nu este necesar ca robotul să se îndepărteze mai întâi de cale.

Comenzile finale includ anticiparea comună:
\begin{equation}
  v_c=\mathrm{sat}(v_r+\delta v,0,v_{\max}),\qquad
  \omega_c=\mathrm{sat}(\omega_r+\delta\omega,
  -\omega_{\max},\omega_{\max}).
\end{equation}
Limitele curente sunt $v_{\max}=1$ m/s și $\omega_{\max}=1{,}5$ rad/s. Același
bloc de anticipare, saturație și siguranță trebuie reutilizat fără modificări
de LQR și MPC. Astfel, diferența dintre experimente va proveni din legea de
feedback, nu din tratamentul diferit al actuatoarelor.

\section{Localizare și evaluare}

Regulatorul primește poziția estimată de EKF din odometria roților și IMU.
Poziția exactă Gazebo nu este folosită în bucla de control. Un evaluator separat
primește ground truth numai pentru calculul performanței. Ceasul evaluatorului
este sincronizat cu primul eșantion al regulatorului, evitând deplasarea
artificială în timp dintre referință și mișcarea măsurată.

Sunt înregistrate două familii de erori:
\begin{itemize}
  \item erorile temporale $e_x$, $e_y$, $e_{\theta}$ și eroarea euclidiană față
  de robotul virtual;
  \item eroarea transversală geometrică și eroarea față de tangenta căii,
  calculate prin proiecția pe cale.
\end{itemize}
Prima familie măsoară urmărirea traiectoriei, iar a doua arată calitatea căii
parcurse independent de întârzierea temporală. Fișierele de sumar includ RMSE,
valorile maxime, eroarea terminală, durata, activitatea normalizată a comenzilor
și eroarea de localizare.

\section{Limitare identificată în Gazebo Classic}

Într-o primă versiune s-a permis o comandă liniară negativă pentru revenirea
după depășirea punctului final. Testul a evidențiat o divergență severă între
EKF și ground truth. Investigația sursei pluginului
\texttt{gazebo\_ros\_diff\_drive} a arătat că, în modul de odometrie
\texttt{ENCODER}, viteza liniară este calculată ca
\begin{equation}
  v_{odom}=\frac{\sqrt{\Delta x^2+\Delta y^2}}{\Delta t},
\end{equation}
deci este întotdeauna pozitivă. Semnul mersului înapoi este pierdut. Sursa
oficială poate fi consultată la
\url{https://github.com/ros-simulation/gazebo_ros_pkgs/blob/ros2/gazebo_plugins/src/gazebo_ros_diff_drive.cpp}.

Folosirea ground truth pentru a corecta această situație ar anula scopul
fuziunii senzoriale. Soluția adoptată este mișcarea liniară înainte, comună
tuturor regulatoarelor, și un profil de frânare suficient de conservator pentru
a evita necesitatea corecției inverse. Această alegere trebuie menționată în
lucrare drept o limitare a infrastructurii de simulare, nu drept o proprietate
generală a regulatorului.

\section{Validare realizată}

Compilarea pachetelor s-a încheiat fără erori. Suita automată conține 99 de
teste trecute, fără eșecuri; 21 de verificări ale analizorului static au fost
omise de instrument deoarece versiunea instalată este cunoscută ca lentă.
Testele includ profilul temporal, independența referinței față de progresul
robotului, transformarea erorilor în reper mobil, cele trei bucle PID, compunerea
celor două comenzi și saturațiile comune.

Tabelul~\ref{tab:smoke} prezintă rulări Gazebo individuale, folosite numai ca
validare funcțională. Ele nu reprezintă încă experimentul statistic final.

\begin{table}[h]
\centering
\caption{Rezultate preliminare cu referință temporală și comandă înainte}
\label{tab:smoke}
\begin{tabular}{lrrrr}
\toprule
Traseu & Durată [s] & Eroare finală reală [m] & RMSE poziție [m] & RMSE CTE [m]\\
\midrule
Dreaptă, 5 m & 14,93 & 0,0578 & 0,0493 & 0,0026\\
Curbă & 20,30 & 0,0561 & 0,0776 & 0,0494\\
Cerc & 25,40 & 0,0290 & 0,0610 & 0,0138\\
\bottomrule
\end{tabular}
\end{table}

Testul pe traseul în formă de opt a fost oprit la cererea utilizatorului înainte
de finalizare și nu este raportat ca rezultat. Valorile mai mari ale erorii
temporale față de vechea eroare geometrică sunt normale: robotul este acum
penalizat și pentru întârzierea sau avansul față de referința temporală.

\section{Fișiere și componente introduse}

Managerul temporal este implementat în
\texttt{trajectory\_reference\_manager.hpp/.cpp}. Politica comună de comandă se
află în \texttt{motion\_command\_policy.hpp/.cpp}. Regulatorul PID cascadat și
nodul ROS au fost extinse pentru $\delta v$ și $\delta\omega$. Evaluatorul de
ground truth folosește aceeași referință temporală și același moment de start.
Scripturile de benchmark au fost adaptate la noile coloane și produc rezumate
compatibile cu importul ulterior în MATLAB.

\section{Pașii următori}

Înaintea comparației finale sunt necesare:
\begin{enumerate}
  \item finalizarea validării pe traseul în opt și pe cazul perturbat;
  \item înlocuirea sau netezirea colțurilor cu discontinuități de tangentă,
  deoarece o schimbare instantanee de $90^{\circ}$ nu este o referință cinematic
  realizabilă la viteză nenulă;
  \item discretizarea modelului liniar variabil în timp la perioada regulatorului;
  \item implementarea LQR cu intrările $\delta v$ și $\delta\omega$;
  \item implementarea MPC pe același model, cu aceleași saturații;
  \item reglarea separată a fiecărui regulator, apoi rulări repetate cu aceleași
  trasee, perturbații și semințe Gazebo;
  \item analiza statistică în MATLAB folosind medie, abatere standard, intervale
  de încredere și compromisul dintre precizie și activitatea comenzilor.
\end{enumerate}

\section{Concluzie intermediară}

Proiectul a trecut efectiv la urmărirea unei traiectorii: referința are poziție,
orientare și două viteze dependente de timp, iar PID controlează atât viteza
liniară, cât și viteza unghiulară. Interfața rezultată este adecvată unei
comparații PID--LQR--MPC pe aceleași variabile. Implementarea este compilată și
testată automat, iar trei probe Gazebo confirmă funcționarea de bază. Rezultatele
de mai sus trebuie tratate drept validare preliminară; concluziile comparative
vor fi formulate numai după implementarea celorlalte regulatoare și efectuarea
matricei experimentale repetate.

\end{document}
